%% kapa: exam=1 round=20 year=2009 period=2 form=A subject=appraisal_law subject_ko=부동산관계법규 seq=24 number=64
%% source: Q-Net 감정평가사 기출문제 (cst00309) artlSeq=5024178 file=extracted/1008/5024178/2009년 제20회 감정평가사 1차시험 기출문제/2교시.hwp
%% conversion: kordoc 4.24.51 -> tools/gen1.py (auto); review=proofread-sonnet; flags=none
\begin{question}{64}
부동산 가격공시 및 감정평가에 관한 법령상 감정평가업자에 대한 과징금의 부과 및 징수에 관한 설명으로 옳지 \uline{않은} 것은?
\begin{choices}
  \item 국토해양부장관은 감정평가업자에게 업무정지처분을 하여야 하는 경우로서 그 업무정지처분이 공익을 해칠 우려가 있는 경우에는 업무정지처분에 갈음하여 과징금을 부과할 수 있다.
  \item 과징금의 최고액은 5천만원(감정평가법인인 경우는 5억원)이며, 과징금의 부과기준은 대통령령으로 정한다.
  \item 합병 전 감정평가법인의 위반행위에 대해서도 합병 후 존속되는 감정평가법인의 행위로 보아 과징금을 부과할 수 있다.
  \item 과징금부과처분에 대하여 이의신청을 제기한 경우 그에 대한 국토해양부장관의 결정에 불복하여 다시 행정심판을 청구할 수는 없다.
  \item 국토해양부장관은 과징금을 부과하는 때에는 그 위반행위의 종별과 해당 과징금의 금액을 명시하여 이를 납부할 것을 서면으로 통지하여야 한다.
\end{choices}
\end{question}
\answer{4}
